\chapter{Sudden infant death syndrome SIDS}
\index{Sudden infant death syndrome SIDS}

\medskip
\noindent\textit{Educational reference; always seek professional advice for individual care.}

\section*{Definition and Overview}
Sudden infant death syndrome, widely known as SIDS, is the sudden, unexpected death of an otherwise healthy baby, usually during sleep, that remains unexplained after a thorough investigation, including a complete autopsy, examination of the scene, and review of the baby's medical history. It is sometimes called "cot death" because it most often occurs while the baby is asleep in a cot. SIDS is one category of a broader group of events called sudden unexpected death in infancy (SUDI), which also includes deaths that are later explained by infection, accidents, or other medical conditions. The distinction matters because the label SIDS is only applied once every other explanation has been excluded.

SIDS is a devastating event for families, and it is a major reason why safe-sleep advice is given such prominence in local maternity and child health care. The great majority of SIDS deaths are associated with modifiable risk factors, and the "Reduce the Risks" campaign has been remarkably successful: local SIDS rates fell by about 80 percent after parents were advised to place babies on their backs to sleep. This chapter explains what SIDS is, who is at risk, and the theories about why it happens. It also gives the current local safe-sleep recommendations in detail, because following them is the single most effective way parents can reduce the risk. The chapter ends with a note of compassion: families who lose a baby to SIDS deserve understanding and support, never blame.

\section*{Epidemiology}
SIDS most commonly occurs in babies between one and six months of age, with a peak between two and four months. It is rare in the first month of life and uncommon after twelve months. Boys are affected slightly more often than girls. Before the safe-sleep campaigns, SIDS was the leading cause of death for infants between one month and one year in many countries and many other countries, with rates around two deaths per thousand live births. Since the early 1990s, the rate in many countries has fallen to well under one per thousand, and SIDS is now far less common, though it has not disappeared.

Certain populations carry a higher risk. Rates are higher among Aboriginal and Torres Strait Islander infants, among babies from families living with disadvantage, and in some geographic regions, reflecting higher rates of maternal smoking and other risk factors. Risk factors include sleeping on the tummy or side, bed-sharing in certain circumstances, exposure to tobacco smoke, overheating, prematurity and low birth weight, and young maternal age. The good news is that these are largely modifiable. some countries were one of the first countries to adopt the back-to-sleep message, and its experience shows clearly that public education saves lives.

\section*{Aetiology and Pathophysiology}
The exact cause of SIDS remains unknown, and it is likely that no single cause exists. The leading scientific model, called the "triple risk" hypothesis, proposes that SIDS occurs when three conditions overlap: a vulnerable infant who is at a critical stage of development, a critical period in the development of sleep and breathing control, and an external stressor such as face-down sleep, overheating, or a respiratory infection. In this model, a baby who would normally arouse from sleep in response to breathing difficulty or low oxygen fails to do so, and the protective reflexes that restore breathing do not fire.

The brainstem, which controls breathing, heart rate, and arousal from sleep, appears to be involved. Studies have found subtle abnormalities in the brainstems of some SIDS babies, in areas involved in serotonin signalling and in responses to oxygen and carbon dioxide. Premature and low-birth-weight babies are more vulnerable because their control systems are immature. Face-down sleeping and soft bedding can trap exhaled air, lowering oxygen and raising carbon dioxide, which should trigger arousal but does not in a vulnerable baby.

\begin{itemize}
    \item \textbf{Stomach sleeping:} the strongest single risk factor; sleeping on the tummy roughly quadruples the risk, and side sleeping is also unsafe
    \item \textbf{Bed-sharing:} sharing a bed with an adult increases risk, particularly in the first three to four months, if the mother smokes, or when the adult is under the influence of alcohol or drugs
    \item \textbf{Maternal smoking:} both smoking during pregnancy and exposing the baby to second-hand smoke increase risk substantially
    \item \textbf{Overheating:} excessive clothing, blankets, and warm room temperature raise the risk
    \item \textbf{Soft bedding:} pillows, quilts, doonas, sheepskins, soft mattresses, and toys in the cot can cover the face and obstruct breathing
    \item \textbf{Prematurity and low birth weight:} immature control of breathing and arousal make these babies more vulnerable
    \item \textbf{Alcohol and drug use by caregivers:} reduce the caregiver's ability to respond to the baby and increase risk during bed-sharing
\end{itemize}

\section*{Clinical Features}
SIDS is, by definition, a death, and it has no symptoms that predict it. However, some features of the sleep environment and the baby's circumstances are associated with higher risk, and it is these that parents and health professionals watch for.

\begin{itemize}
    \item \textbf{Occurrence during sleep:} most SIDS deaths happen during sleep, usually at night, and many babies are found face down or with bedding over their faces
    \item \textbf{Age at peak risk:} most deaths occur between one and six months of age
    \item \textbf{An apparently healthy baby:} the baby usually appears well beforehand, with no serious illness
    \item \textbf{Recent mild illness:} a minor cold or fever is sometimes present in the days before death
    \item \textbf{Co-sleeping history:} a significant proportion of deaths occur while the baby is sharing a bed, sofa, or chair with an adult
    \item \textbf{Exposure to cigarette smoke:} maternal smoking in pregnancy or smoke in the home is common in SIDS histories
    \item \textbf{Unsafe sleep surfaces:} soft mattresses, sofa sleeping, or unsafe cots are frequently found at the scene
\end{itemize}

\section*{Differential Diagnosis}
Because SIDS is a diagnosis of exclusion, clinicians must consider other causes of sudden infant death before applying the label.

\begin{itemize}
    \item \textbf{Accidental suffocation:} death from overlay (an adult rolling onto the baby), entrapment between mattress and wall, or covering of the airway by bedding or a soft surface
    \item \textbf{Infection:} overwhelming bacterial infection, such as from some strains of staphylococcus or meningococcus, or severe viral infection, can cause rapid death
    \item \textbf{Metabolic disorders:} rare inherited conditions such as medium-chain acyl-CoA dehydrogenase (MCAD) deficiency cause sudden death, usually detected by newborn screening
    \item \textbf{Cardiac conditions:} inherited heart rhythm disorders such as long QT syndrome can cause sudden death in infancy
    \item \textbf{Neurological conditions:} severe epilepsy and some brain malformations can present as sudden death
    \item \textbf{Non-accidental injury:} inflicted head injury or suffocation must be considered; thorough investigation by the coroner and child-protection authorities distinguishes this from SIDS
    \item \textbf{Intoxication:} poisoning from medicines or household substances is a rare cause of unexpected infant death
\end{itemize}

\section*{When to Seek Medical Care}
There are no warning signs that reliably predict SIDS, so this section is about practical vigilance rather than symptoms to watch for. A baby who is unusually floppy, difficult to rouse, breathing rapidly or with difficulty, or who has a high fever and is not alert should be seen by a doctor promptly. Feeding, breathing, and behaviour changes matter.

\textbf{Contact your local emergency services for:}
\begin{itemize}
    \item A baby who is not breathing, unresponsive, or blue or very pale
    \item A seizure or a baby who cannot be woken
    \item Severe breathing difficulty, with the chest sucking in or the baby making grunting sounds
    \item A baby who is limp, floppy, or unconscious
\end{itemize}

If the unthinkable happens and a baby is found not breathing, begin CPR immediately and have someone contact local emergency services. Do not delay. Bystander CPR before emergency services arrive can be life-saving, and parents are encouraged to learn infant CPR.

\section*{Diagnosis and Investigations}
The diagnosis of SIDS is made only after a full investigation has ruled out other causes.

\begin{itemize}
    \item \textbf{Full autopsy:} a thorough examination by a forensic pathologist, including detailed examination of the organs and brain
    \item \textbf{Scene investigation:} trained investigators examine the place where the baby died, including the sleep surface and circumstances
    \item \textbf{Review of medical history:} the baby's health record, feeding history, and any recent illness are reviewed
    \item \textbf{Microbiology and toxicology:} tests for infection and for any medicines or substances
    \item \textbf{Metabolic and genetic testing:} newborn screening results and specialised testing look for inherited metabolic or cardiac conditions
    \item \textbf{Radiology:} skeletal X-rays are performed to look for injury or unsuspected bone disease
    \item \textbf{Supportive interview with the family:} conducted sensitively, because the family's information is essential to the investigation
\end{itemize}

\section*{Management}
SIDS is a death, not a treatable condition, so the "management" is the care of the family and the prevention of further deaths.

\begin{itemize}
    \item \textbf{Bereavement support:} families who lose a baby need immediate and ongoing grief support from professionals who understand SIDS
    \item \textbf{Coronial process support:} families are guided through the coronial investigation, which can feel overwhelming and prolonged
    \item \textbf{Counselling:} grief counselling and support groups connect bereaved parents with others who have shared the experience
    \item \textbf{Genetic counselling:} where inherited conditions are identified, families receive information about risks for future children
    \item \textbf{Care for subsequent babies:} families who go on to have another baby receive enhanced monitoring and repeated safe-sleep education
    \item \textbf{Peer support organisations:} bodies such as Red Nose provide resources, counselling, and connection with other bereaved families
    \item \textbf{Staff debriefing:} health and emergency staff affected by the death receive support
\end{itemize}

\section*{Complications}
\begin{itemize}
    \item \textbf{Intense parental grief:} the sudden, unexplained death of a baby is among the most traumatic of human experiences, and grief can be complicated by guilt
    \item \textbf{Guilt and self-blame:} parents often blame themselves even when there was nothing they could have done
    \item \textbf{Relationship strain:} the loss can put enormous pressure on the parents' relationship and on siblings
    \item \textbf{Anxiety in future parenting:} parents who have lost a baby often experience severe anxiety during subsequent pregnancies and parenting
    \item \textbf{Mental health difficulties:} depression, post-traumatic stress, and prolonged complicated grief are common after the loss
    \item \textbf{Increased risk in subsequent pregnancies:} families who have lost one baby to SIDS may face an elevated risk with future babies, which is why careful antenatal and infant follow-up is offered
\end{itemize}

\section*{Prognosis and Outcomes}
SIDS itself has no survivors, and the term "prognosis" here applies to prevention and to the family's journey afterwards. The epidemiology is hopeful in one important way: most SIDS deaths are linked to modifiable risk factors, and following safe-sleep advice reduces the risk dramatically. The fall in local SIDS rates since the back-to-sleep campaigns is one of the great public health successes of modern child health, and it demonstrates that the risk is substantially within families' control, supported by health professionals. For the families left behind, outcomes vary widely. Many parents, with time, professional support, and connection to other bereaved families, find ways to carry their grief and rebuild their lives, and many go on to have healthy subsequent children. Red Nose and similar organisations provide the specialised grief support that helps make this possible.

\section*{Prevention and Self-Care}
\begin{itemize}
    \item Always place your baby on their back to sleep, for every sleep, day and night
    \item Use a firm, flat mattress in a safe cot that meets local standards, with no pillows, quilts, doonas, sheepskins, bumpers, or soft toys
    \item Keep your baby's head and face uncovered during sleep, with bedding tucked in firmly at chest level
    \item Sleep your baby in the same room as you, in their own cot or bassinet, for at least the first six to twelve months
    \item Avoid bed-sharing, particularly in the first four months, and never share a bed if the baby was born prematurely, if the mother smoked in pregnancy, or if anyone has consumed alcohol, drugs, or medicines that cause drowsiness
    \item Never sleep with your baby on a sofa or armchair, which is especially dangerous
    \item Do not smoke during pregnancy and keep the home and car smoke-free
    \item Keep the room at a comfortable temperature and dress the baby for the room, not the season; check for signs of overheating such as sweating or a hot chest
    \item Breastfeed if you can, offer a dummy (pacifier) for sleeps once breastfeeding is established, and keep your baby immunised, as these are each associated with a lower risk
    \item Learn infant CPR so you are prepared in an emergency
\end{itemize}

\section*{Sources and Further Reading}
The following organisations provide reliable information about SIDS and safe infant sleeping.

\begin{itemize}
    \item Red Nose (Australia). \textit{Safe sleeping advice for babies and families.} https://rednose.org.au/
    \item healthdirect Australia. \textit{SIDS and safe infant sleeping.} https://www.healthdirect.gov.au/
    \item NHS. \textit{Sudden infant death syndrome: reducing the risk.} https://www.nhs.uk/
    \item Raising Children Network. \textit{Safe sleeping for babies.} https://raisingchildren.net.au/
    \item Australian Department of Health and Aged Care. \textit{Safe sleeping and child health resources.} https://www.health.gov.au/
    \item World Health Organization (WHO). \textit{Child health and safe sleep guidance.} https://www.who.int/
\end{itemize}

\clearpage
